Este trabalho propõe um estudo experimental para avaliar comparativamente
alternativas ao MinIO para armazenamento de objetos em arquiteturas de
\textit{Data Lake} híbridas. A pesquisa utilizará o MinIO como \textit{baseline}
e incluirá obrigatoriamente o RustFS, além de outras soluções selecionadas a
partir de um mapeamento sistemático da literatura. Serão analisados desempenho,
compatibilidade com a API S3, interoperabilidade com AWS, Azure e GCP, consumo de
recursos, custos, licenciamento, maturidade, governança, sustentabilidade e riscos
tecnológicos. O mapeamento sistemático identificará dimensões, métricas e lacunas
relevantes para orientar a construção de um ambiente experimental padronizado.

\begin{keywords}
Armazenamento de objetos, Data Lake, MinIO, RustFS, API S3, nuvem híbrida.
\end{keywords}
